\documentclass[11pt,a4paper]{article}
\usepackage[T1]{fontenc}
\usepackage[utf8]{inputenc}
\usepackage[spanish,es-noquoting,es-nodecimaldot]{babel}
\usepackage{lmodern}
\usepackage[margin=2.5cm]{geometry}
\usepackage{amsmath,amssymb,amsthm}
\usepackage{booktabs}
\usepackage{microtype}
\usepackage{needspace}
\usepackage[hidelinks]{hyperref}
\hypersetup{pdftitle={Forgetting con reemplazo aleatorio: dos pruebas de la política de intervalo},pdfauthor={Notas de investigación}}
\newtheorem{proposicion}{Proposición}
\renewcommand{\proofname}{Demostración}
\setlength{\parindent}{0pt}
\setlength{\parskip}{0.45em}
\setlength{\emergencystretch}{1.5em}
\allowdisplaybreaks[1]
\raggedbottom
\title{Forgetting con reemplazo aleatorio\\[0.35em]
\large Dos pruebas de la política de intervalo}
\author{Progress notes}
\date{3 de noviembre de 2026}

\begin{document}
\maketitle

Esta nota explica por qué, cuando la memoria está llena, el vendedor
elige el producto de mayor calidad en un solo intervalo de estados.
La primera prueba sigue las ecuaciones de Bellman. La segunda usa
programación lineal para demostrar que la mejor política de dos
umbrales es óptima en el problema completo.

Las dos pruebas comparten un paso: demostrar que la ventaja de elegir
el producto 2 no puede ser positiva en dos regiones separadas por una
región donde es no positiva. La prueba por programación lineal utiliza
ese paso y lo combina con un certificado de optimalidad.

\section{El modelo y el resultado que queremos probar}

\subsection{Estado, productos y demanda}

La capacidad $n$ es un entero positivo. La memoria contiene exactamente
$n$ resultados. El estado $s$ es el
número de éxitos recordados, de modo que $s\in\{0,\ldots,n\}$.
El vendedor puede entregar dos productos. El producto 2 tiene mayor
probabilidad de éxito y también mayor coste. Suponemos
\begin{equation}
0<p_0<1,\qquad 0<p_1<p_2<1,\qquad
0\le c_1<c_2<R,\qquad 0<\gamma<1.
\label{eq:supuestos}
\end{equation}
Aquí $p_0$ es la calidad de la opción externa, $p_i$ es la probabilidad
de éxito del producto $i$, $c_i$ es su coste, $R$ es el ingreso por
venta y $\gamma$ es el factor de descuento por período calendario.

Con prior Beta$(1,1)$ y la regla de elección del modelo, la probabilidad
de comprar al vendedor en el estado $s$ es
\begin{equation}
d_s=\Pr\{\operatorname{Beta}(s+1,n-s+1)>p_0\}
   =\sum_{j=0}^{s}\binom{n+1}{j}p_0^j(1-p_0)^{n+1-j}.
\label{eq:demand}
\end{equation}
La segunda expresión es una identidad entre la distribución Beta y la
binomial. Nos permitirá estudiar cómo cambia la demanda al recordar
un éxito adicional. En particular, $d_s>0$ para todos los estados.

Después de cada compra, el resultado nuevo reemplaza una observación
antigua elegida uniformemente, con independencia de si el resultado
nuevo fue un éxito o un fracaso. Por ello, condicional a una compra
y a usar el producto $i$,
\begin{center}
\begin{tabular}{lll}
\toprule
Estado siguiente & Probabilidad & Qué ocurrió \\
\midrule
$s+1$ & $p_i(1-s/n)$ & Éxito nuevo; se elimina un fracaso. \\
$s-1$ & $(1-p_i)s/n$ & Fracaso nuevo; se elimina un éxito. \\
$s$ & Probabilidad restante & El dato nuevo y el eliminado coinciden. \\
\bottomrule
\end{tabular}
\end{center}
En $s=0$ no se puede bajar y en $s=n$ no se puede subir. Si el
consumidor elige la opción externa, el estado no cambia y el beneficio
del período es cero.

\subsection{Qué significa una política de intervalo}

Una política estacionaria asigna un producto a cada estado, sin
depender del período calendario. Queremos demostrar que existe una
política óptima de la forma
\begin{equation}
\pi(s)=
\begin{cases}
2,&L\le s\le U,\\
1,&\text{en los demás estados}.
\end{cases}
\label{eq:band-policy}
\end{equation}
Es decir, al recorrer los estados de menor a mayor, la política tiene
el patrón
\[
\underbrace{1,\ldots,1}_{s<L}\ ,\qquad
\underbrace{2,\ldots,2}_{L\le s\le U}\ ,\qquad
\underbrace{1,\ldots,1}_{s>U}.
\]
Se permiten intervalos que llegan a un extremo, un único estado,
todo el espacio de estados o ningún estado. Por tanto, el resultado
incluye políticas con un solo umbral y políticas constantes.

\begin{proposicion}\label{prop:intervalo}
Bajo los supuestos \eqref{eq:supuestos}, con prior Beta$(1,1)$ y
memoria llena, existe una política óptima de intervalo. Podemos
elegir el producto 1 siempre que ambas acciones den el mismo valor.
Si no hay empates, la política óptima estacionaria es única.
\end{proposicion}

El resultado no requiere restricciones adicionales sobre los parámetros
de \eqref{eq:supuestos}. Se refiere a memoria llena; al final de la
nota explicamos por qué el llenado inicial debe tratarse por separado.

\section{Primera prueba: estudiar la ventaja del producto 2}
\label{sec:primera}

\subsection{Qué comparación determina la acción}

Sea $V_s$ el beneficio descontado óptimo a partir del estado $s$.
La ecuación de Bellman considera primero si ocurre una compra y,
si ocurre, compara los dos productos. Usando las probabilidades de
transición de la tabla y agrupando los términos en $V_s$, obtenemos
\begin{equation}
\begin{aligned}
(1-\gamma)V_s=d_s\max_{i\in\{1,2\}}\biggl\{
&R-c_i+\gamma p_i(1-s/n)(V_{s+1}-V_s)\\
&-\gamma(1-p_i)(s/n)(V_s-V_{s-1})\biggr\}.
\end{aligned}
\label{eq:bellman}
\end{equation}
En los extremos omitimos el término cuyo coeficiente es cero. Esta
convención se mantiene en todas las expresiones con estados vecinos.

Para comparar acciones, escribimos únicamente las dos diferencias
de parámetros que aparecen una y otra vez:
\[
\Delta p=p_2-p_1,\qquad \Delta c=c_2-c_1.
\]
Si usamos $V$ como valor de continuación, la ventaja del producto 2
frente al producto 1, condicional a una compra, es
\begin{equation}
G_s=-\Delta c+\gamma\Delta p
\left[(1-s/n)(V_{s+1}-V_s)+(s/n)(V_s-V_{s-1})\right].
\label{eq:advantage}
\end{equation}
El primer término es el coste adicional. El segundo es el beneficio
futuro de aumentar la probabilidad de éxito. Los dos incrementos del
valor aparecen porque un éxito nuevo puede elevar el estado, o evitar
la caída que habría producido un fracaso.

El producto 2 es estrictamente mejor cuando $G_s>0$. Así, basta
demostrar que los estados con ventaja positiva forman un intervalo.
Para un problema finito y descontado como este existe una política
óptima estacionaria: podemos elegir en cada estado una acción que
alcance el máximo de Bellman.

\subsection{Una transformación que cancela el coste adicional}

La constante
\[
k=\frac{\Delta c}{\gamma\Delta p}
\]
es el incremento de valor que, después de descontar y multiplicar por
la mejora en la probabilidad de éxito, compensa exactamente el coste
adicional. Restemos ese incremento de referencia por cada éxito
recordado:
\begin{equation}
W_s=V_s-ks.
\label{eq:transform}
\end{equation}
Por ejemplo, $V_{s+1}-V_s=(W_{s+1}-W_s)+k$. Al sustituirlo en
\eqref{eq:advantage}, los dos pesos suman uno y el término
$\gamma\Delta p\,k$ cancela $\Delta c$. Queda
\begin{equation}
G_s=\gamma\Delta p
\left[(1-s/n)(W_{s+1}-W_s)+(s/n)(W_s-W_{s-1})\right].
\label{eq:advantage-transformed}
\end{equation}
Ahora solo necesitamos estudiar los signos de los incrementos de $W$.

La misma sustitución en Bellman simplifica también el beneficio.
Para cada producto se cumple
\[
R-c_i+\gamma k p_i=R-c_1+\gamma k p_1.
\]
Por eso, el término que queda después de restar $ks$ es el mismo
para ambos productos. Lo llamaremos $F_s$:
\begin{equation}
F_s=d_s\left[R-c_1+\gamma k(p_1-s/n)\right]
       -(1-\gamma)ks.
\label{eq:transformed-reward}
\end{equation}
La ecuación para $W$ queda
\begin{equation}
\begin{aligned}
(1-\gamma)W_s=F_s+\gamma d_s\max_{i\in\{1,2\}}\biggl\{
&p_i(1-s/n)(W_{s+1}-W_s)\\
&-(1-p_i)(s/n)(W_s-W_{s-1})\biggr\}.
\end{aligned}
\label{eq:transformed-bellman}
\end{equation}
Este es un cambio algebraico exacto. La cantidad $F_s$ es una ayuda
para la prueba, no un nuevo beneficio económico que se añada al modelo.

\subsection{Dónde son positivos los incrementos de F}
\label{sec:signos-F}

Usaremos $\Delta F_j=F_j-F_{j-1}$ para el cambio de $F$ entre estados
consecutivos. Del mismo modo, por \eqref{eq:demand},
\begin{equation}
\Delta d_j=d_j-d_{j-1}
 =\binom{n+1}{j}p_0^j(1-p_0)^{n+1-j}>0,
\qquad j=1,\ldots,n.
\label{eq:delta-demand}
\end{equation}

Para que las expresiones de este paso sean legibles, usaremos
temporalmente una sola constante:
\[
\tau=n\left[p_1+\frac{R-c_1}{\gamma k}\right]>0.
\]
Con ella, \eqref{eq:transformed-reward} y su diferencia se escriben
\begin{align}
F_s&=\frac{\gamma k}{n}d_s(\tau-s)-(1-\gamma)ks,
\notag\\
\Delta F_j&=\frac{\gamma k}{n}
 \left[\Delta d_j(\tau-j)-d_{j-1}\right]-(1-\gamma)k.
\label{eq:delta-reward}
\end{align}
Si $j\ge\tau$, todos los términos de la última expresión son no
positivos y el último es estrictamente negativo. Por tanto,
$\Delta F_j<0$. Solo necesitamos estudiar los índices $j<\tau$.

La razón entre dos probabilidades binomiales consecutivas es
\[
\frac{\Delta d_{j+1}}{\Delta d_j}
  =\frac{n+1-j}{j+1}\frac{p_0}{1-p_0}.
\]
Al restar dos expresiones consecutivas de \eqref{eq:delta-reward}
y usar esa razón, obtenemos, para $j=1,\ldots,n-1$,
\begin{equation}
\begin{aligned}
\Delta F_{j+1}-\Delta F_j
={}&\frac{\gamma k\,\Delta d_j}{n(j+1)(1-p_0)}\\[-0.2em]
&\times\left[(\tau-j-1)((n+2)p_0-j-1)
                     -2(j+1)(1-p_0)\right].
\end{aligned}
\label{eq:second-difference}
\end{equation}
El factor fuera del corchete es positivo. Para entender el signo del
corchete, llamemos $x=j+1$ y observemos la expresión
\begin{equation}
(\tau-x)((n+2)p_0-x)-2(1-p_0)x.
\label{eq:sign-expression}
\end{equation}
Mientras $0<x<\min\{\tau,(n+2)p_0\}$, los dos factores del producto
son positivos y disminuyen con $x$. Su producto disminuye, mientras
el término que restamos aumenta. La expresión empieza positiva en
$x=0$ y es negativa al llegar al menor de esos dos límites. Por eso
su signo solo puede pasar una vez de positivo a negativo en ese tramo.
Si $(n+2)p_0\le x<\tau$, el producto ya es no positivo y el término
que restamos es positivo, así que la expresión sigue siendo negativa.

Esto demuestra que, entre los índices donde $\Delta F_j$ puede ser
positivo, sus valores primero aumentan y después disminuyen. A partir
de $j\ge\tau$, ya sabemos que son negativos. En consecuencia,
\begin{equation}
\{j:\Delta F_j>0\}\quad\text{es un intervalo, posiblemente vacío.}
\label{eq:positive-F}
\end{equation}
En palabras: al recorrer los incrementos de $F$, no podemos encontrar
uno positivo, luego uno no positivo y después otro positivo.
Esta propiedad del signo es todo lo que necesitamos; no hace falta
demostrar que toda la función $F$ sea cóncava.

\subsection{La misma propiedad pasa a los incrementos del valor}
\label{sec:signos-W}

En este paso fijamos \emph{cualquier} política estacionaria $\pi$.
Escribimos $V^\pi$ para su valor y $W_s^\pi=V_s^\pi-ks$ para su
valor transformado. No suponemos que la política sea óptima ni que
tenga forma de intervalo. Esta generalidad servirá también en la
prueba por programación lineal.

Para aligerar únicamente las ecuaciones de este paso, escribimos
\[
\lambda_s=\gamma d_s p_{\pi(s)}(1-s/n),\qquad
\mu_s=\gamma d_s(1-p_{\pi(s)})s/n.
\]
Son las probabilidades de subir y bajar, respectivamente, multiplicadas
por el descuento. Ambas son no negativas; $\lambda_n=\mu_0=0$.
La evaluación de la política satisface
\begin{equation}
(1-\gamma)W_s^\pi=F_s
 +\lambda_s(W_{s+1}^\pi-W_s^\pi)
 -\mu_s(W_s^\pi-W_{s-1}^\pi).
\label{eq:fixed-policy}
\end{equation}

Llamemos $w_j=W_j^\pi-W_{j-1}^\pi$ a los incrementos que queremos
estudiar. Restar la ecuación de $j-1$ de la ecuación de $j$ da
\begin{equation}
\bigl[(1-\gamma)+\mu_j+\lambda_{j-1}\bigr]w_j
 -\mu_{j-1}w_{j-1}-\lambda_jw_{j+1}
 =\Delta F_j.
\label{eq:tridiagonal}
\end{equation}
Cada ecuación conecta un incremento solo con sus dos vecinos.
Por eso el sistema es \emph{tridiagonal}: su matriz tiene valores
distintos de cero únicamente en la diagonal y en las dos diagonales
adyacentes.

Resolver este sistema no puede convertir un único intervalo positivo
del lado derecho en dos intervalos positivos de la solución. Podemos
verlo con la eliminación de Gauss, sin recurrir a un resultado
adicional sobre matrices.

Al eliminar sucesivamente las incógnitas de izquierda a derecha,
los coeficientes de la diagonal se actualizan así:
\begin{align}
a_1&=(1-\gamma)+\mu_1+\lambda_0,\notag\\
a_j&=(1-\gamma)+\mu_j+\lambda_{j-1}
       -\frac{\mu_{j-1}\lambda_{j-1}}{a_{j-1}},
\qquad j=2,\ldots,n.
\label{eq:pivots}
\end{align}
Estos coeficientes son positivos. En efecto, $a_1>\mu_1$. Si
$a_{j-1}>\mu_{j-1}$, entonces
\[
a_j=(1-\gamma)+\mu_j+
\lambda_{j-1}\left(1-\frac{\mu_{j-1}}{a_{j-1}}\right)>\mu_j.
\]
Así, podemos hacer todas las divisiones de la eliminación.

Los lados derechos actualizados, que llamamos $y_j$, y la sustitución
final son
\begin{align}
y_1&=\Delta F_1,
&y_j&=\Delta F_j+\frac{\mu_{j-1}}{a_{j-1}}y_{j-1},
\label{eq:forward}\\
w_n&=\frac{y_n}{a_n},
&w_j&=\frac{y_j+\lambda_jw_{j+1}}{a_j}.
\label{eq:backward}
\end{align}
La primera línea se recorre hacia la derecha y la segunda hacia la
izquierda. Todos los pesos que multiplican los valores ya calculados
son no negativos.

Para entender qué ocurre con los signos en \eqref{eq:forward},
recorramos los índices de izquierda a derecha:
\begin{itemize}
\item Antes del intervalo donde $\Delta F_j>0$, todos los valores
calculados son no positivos.
\item Dentro de ese intervalo, una vez que $y_j$ se hace positivo,
los siguientes siguen siendo positivos: sumamos una entrada positiva
y un múltiplo no negativo de un valor positivo.
\item Después del intervalo, las entradas son no positivas.
Una vez que $y_j$ se hace no positivo, ya no puede volver a ser positivo.
\end{itemize}
Por tanto, los valores positivos de $y_j$ forman un intervalo. El
mismo razonamiento, leyendo de derecha a izquierda y dividiendo por
los números positivos $a_j$, se aplica a \eqref{eq:backward}.
Concluimos que
\begin{equation}
\{j:W_j^\pi-W_{j-1}^\pi>0\}
\quad\text{es un intervalo, posiblemente vacío.}
\label{eq:positive-W}
\end{equation}

\subsection{De los incrementos a la política óptima}
\label{sec:greedy-interval}

Sea $G_s^\pi$ la ventaja de elegir producto 2 una vez y después
continuar con la política $\pi$. Se obtiene sustituyendo $V^\pi$
por $V$ en \eqref{eq:advantage}. Por
\eqref{eq:advantage-transformed}, su signo es el del promedio
\[
(1-s/n)w_{s+1}+(s/n)w_s.
\]
En los extremos el promedio es $w_1$ o $w_n$.

Si todos los $w_j$ son no positivos, ninguna ventaja es positiva.
En caso contrario, supongamos que los positivos están exactamente
entre $a$ y $b$. Para los estados $s=a,\ldots,b-1$, el promedio
usa dos incrementos positivos, por lo que es positivo. Fuera de los
estados $a-1,\ldots,b$, usa solo incrementos no positivos. Únicamente
los dos estados $a-1$ y $b$ pueden tener cualquiera de los dos signos.
Agregar o quitar estos extremos no puede dividir el intervalo.
Si $a=b$, los únicos candidatos son los dos estados consecutivos
$a-1$ y $a$, de modo que la conclusión también vale.

Hemos demostrado algo útil por sí mismo:
\begin{quote}
Para cualquier política estacionaria fija, los estados en que es
estrictamente mejor elegir producto 2 y después continuar con esa
política forman un intervalo.
\end{quote}

Para terminar la primera prueba, tomemos una política óptima
estacionaria. Su valor es $V$. Elegir producto 2 cuando $G_s>0$ y
producto 1 cuando $G_s\le0$ alcanza el máximo de Bellman en cada
estado, incluidos los empates. Esa elección es óptima y, por lo que
acabamos de demostrar, tiene forma de intervalo. Esto prueba la
Proposición~\ref{prop:intervalo}.\hfill$\square$

\section{Segunda prueba: programación lineal y mejora de políticas}
\label{sec:lp}

La segunda prueba organiza el mismo resultado como un problema de
optimización. Primero escribimos un programa lineal que representa todas
las decisiones posibles. Después mostramos que la mejor política de
intervalo también resuelve ese programa completo.

Necesitaremos una conclusión de la primera prueba: para cualquier
política fija, elegir en cada estado la mejor acción usando su valor de
continuación produce una política de intervalo. Por tanto, esta segunda
prueba comparte con la primera el argumento sobre los signos; la
programación lineal permite convertirlo en un certificado de optimalidad.

\subsection{Escribir Bellman como un programa lineal}

Para incluir en una sola expresión tanto las compras como los períodos
sin compra, escribimos $P_i(s,t)$ para la probabilidad de pasar de $s$ a
$t$ durante un período calendario si se usa el producto $i$. En nuestro
modelo,
\begin{align}
P_i(s,s+1)&=d_s p_i(1-s/n),\notag\\
P_i(s,s-1)&=d_s(1-p_i)s/n,\label{eq:lp-transition}\\
P_i(s,s)&=1-P_i(s,s+1)-P_i(s,s-1).\notag
\end{align}
En los extremos se omiten las transiciones que salen de
$\{0,\ldots,n\}$. Las demás probabilidades son cero. Así, $P_i$
ya incluye el caso en que el consumidor elige la opción externa.

El beneficio esperado de ese período es $d_s(R-c_i)$. El programa lineal
de Bellman es
\begin{equation}
\begin{aligned}
\min_{V_0,\ldots,V_n\in\mathbb R}\quad
    &\sum_{s=0}^{n}V_s\\
\text{sujeto a}\quad
    &V_s\ge d_s(R-c_i)+\gamma\sum_{t=0}^{n}P_i(s,t)V_t,\\
    &\hspace{4cm}s=0,\ldots,n,\quad i=1,2.
\end{aligned}
\label{eq:lp-primal}
\end{equation}

Aunque queremos maximizar los beneficios, aquí minimizamos. La razón es
que cada restricción exige que $V_s$ sea al menos tan grande como el
beneficio de elegir una acción y continuar con los valores $V_t$.
Un vector factible es, por tanto, una cota superior del valor que puede
alcanzar cualquier política. Buscamos la cota más ajustada.

Para verificar esa interpretación, apliquemos una restricción a la
acción que utiliza una política cualquiera. Repitiendo la desigualdad
durante $N$ períodos, $V_s$ domina sus beneficios descontados hasta
$N-1$ más el término de continuación
$\gamma^N\mathbb E[V_{S_N}]$, donde $S_N$ es el estado después de
$N$ períodos. Este último término tiende a cero porque hay
finitos estados y $\gamma<1$. Por tanto, todo vector factible domina
el valor de cualquier política. El valor óptimo satisface las
restricciones y alcanza igualdad con la acción elegida. Es entonces
la solución de \eqref{eq:lp-primal}. La formulación general se encuentra
en \cite{ying-zhu2020}.

\subsection{Qué significa el dual y cómo certifica una política}

El dual de \eqref{eq:lp-primal} introduce una variable $x_{s,i}$ para
cada estado y producto:
\begin{equation}
\begin{aligned}
\max_{x_{s,i}\ge0}\quad
    &\sum_{s=0}^{n}\sum_{i=1}^{2}d_s(R-c_i)x_{s,i}\\
\text{sujeto a}\quad
    &\sum_{i=1}^{2}x_{s,i}
      -\gamma\sum_{t=0}^{n}\sum_{i=1}^{2}P_i(t,s)x_{t,i}=1,\\
    &\hspace{6cm}s=0,\ldots,n.
\end{aligned}
\label{eq:lp-dual}
\end{equation}

Las variables $x_{s,i}$ cuentan visitas descontadas al estado $s$ en las
que la política prescribe el producto $i$ si hay una compra. El lado
derecho es uno porque en el primal
sumamos el valor de todos los estados con el mismo peso. Una forma de
interpretarlo es iniciar una trayectoria en cada estado y sumar sus
visitas descontadas. La restricción expresa un balance: las visitas
totales a $s$ son la visita inicial más las que llegan desde el período
anterior, ponderadas por $\gamma$. Sumando los balances, el número total
de visitas descontadas es $(n+1)/(1-\gamma)$, de modo que es finito.

Ahora tomemos una política candidata $\pi$, por ejemplo, producto 2
entre dos umbrales y producto 1 fuera. Su valor se calcula resolviendo
las ecuaciones lineales
\begin{equation}
V_s^\pi=d_s(R-c_{\pi(s)})
       +\gamma\sum_{t=0}^{n}P_{\pi(s)}(s,t)V_t^\pi.
\label{eq:lp-evaluate}
\end{equation}
Por la estructura de las transiciones, cada ecuación solo involucra
el estado actual y sus vecinos.

Las visitas descontadas generadas por esta política son factibles en
el dual. Además, $x_{s,i}=0$ si $i\ne\pi(s)$. Multiplicar las
ecuaciones \eqref{eq:lp-evaluate} por las visitas correspondientes y
usar los balances del dual da
\begin{equation}
\sum_{s=0}^{n}V_s^\pi
 =\sum_{s=0}^{n}\sum_{i=1}^{2}d_s(R-c_i)x_{s,i}.
\label{eq:lp-equal-objectives}
\end{equation}

Así, para certificar la política solo falta comprobar que $V^\pi$
también satisface las restricciones de la acción alternativa en
\eqref{eq:lp-primal}. Si las satisface, tenemos una solución factible
primal y otra dual con el mismo objetivo. La dualidad débil garantiza
que ambas son óptimas. En términos de \emph{holgura complementaria},
la política solo asigna visitas a acciones cuyas restricciones se
cumplen con igualdad.

Para una candidata que usa producto 2 en $[L,U]$, esta comprobación
consiste en verificar que la ventaja $G_s^\pi$ es no negativa dentro
del intervalo y no positiva fuera. Hay que verificar todos los
estados, no solo los dos extremos.

\subsection{Por qué algún intervalo satisface todas las restricciones}

Hasta aquí sabemos cómo verificar una política candidata. Falta
demostrar que podemos encontrar una candidata óptima dentro de la
familia de intervalos.

Para ello usamos una \emph{mejora de políticas}: dado el valor
$V^\pi$ de una política, elegimos en cada estado la acción con mayor
beneficio actual más valor de continuación. Esta acción se llama
\emph{greedy respecto a $V^\pi$}. No es una decisión basada únicamente
en el beneficio inmediato: incorpora la continuación de la política
que acabamos de evaluar.

La mejora elige producto 2 exactamente cuando $G_s^\pi>0$, y producto
1 en los demás casos. La primera prueba muestra que ese conjunto es
un intervalo para \emph{cualquier} política fija $\pi$. Sus pasos
algebraicos no exigían que la política inicial fuese óptima.

Elijamos ahora la política de intervalo $\bar\pi$ que maximiza
\[
\sum_{s=0}^{n}V_s^\pi
\]
entre todos los intervalos, incluyendo el vacío. Esa elección existe
porque hay un número finito de intervalos. Todavía no suponemos que
sea la mejor entre las políticas de cualquier forma.

Supongamos que $V^{\bar\pi}$ viola alguna restricción del primal.
Construyamos su mejora greedy, $\widehat\pi$. En cada estado, esta
mejora puede elegir la acción que usaba $\bar\pi$, así que
\begin{equation}
d_s(R-c_{\widehat\pi(s)})
 +\gamma\sum_{t=0}^{n}P_{\widehat\pi(s)}(s,t)V_t^{\bar\pi}
 \ge V_s^{\bar\pi}.
\label{eq:lp-one-step}
\end{equation}
La desigualdad es estricta en al menos un estado, precisamente porque
alguna restricción fue violada.

Para ver que la mejora también aumenta el valor total, empecemos con
$V^{(0)}=V^{\bar\pi}$ y evaluemos repetidamente la nueva política:
\[
V_s^{(\ell+1)}=d_s(R-c_{\widehat\pi(s)})
  +\gamma\sum_{t=0}^{n}P_{\widehat\pi(s)}(s,t)V_t^{(\ell)},
\qquad \ell=0,1,\ldots.
\]
Por \eqref{eq:lp-one-step}, el primer paso aumenta o mantiene cada
componente y aumenta estrictamente alguna. Como las probabilidades
son no negativas, el siguiente paso tampoco puede reducir ninguna
componente; lo mismo ocurre en todos los pasos posteriores. Esta
evaluación converge a $V^{\widehat\pi}$: en cada iteración, el mayor
error respecto a ese valor se multiplica, como máximo, por
$\gamma<1$. En consecuencia,
\[
V_s^{\widehat\pi}\ge V_s^{\bar\pi}\quad\text{para todo }s,
\qquad
\sum_{s=0}^{n}V_s^{\widehat\pi}
   >\sum_{s=0}^{n}V_s^{\bar\pi}.
\]

Pero la nueva política $\widehat\pi$ también es un intervalo, por la
conclusión de la primera prueba. Esto contradice que $\bar\pi$ fuera
el mejor intervalo.

Por tanto, $V^{\bar\pi}$ satisface todas las restricciones del primal.
Sus visitas descontadas proporcionan una solución dual con el mismo
objetivo, como en \eqref{eq:lp-equal-objectives}. La mejor política de
intervalo es entonces óptima en el problema completo.\hfill$\square$

El argumento también sugiere un procedimiento: empezar con una
política cualquiera, evaluarla y mejorarla. Después de la primera
mejora, la política ya es un intervalo. Las mejoras posteriores
mantienen esa forma y se detienen cuando se cumplen todas las
restricciones del primal. La propiedad de intervalo utilizada aquí
se refiere a memoria llena; no se extiende automáticamente al llenado
inicial.


\clearpage
\section{Aplicación y alcance del resultado}

Para la calibración
\[
\begin{gathered}
n=20,\quad p_0=0.5,\quad p_1=0.35,\quad p_2=0.8,\\
R=1,\quad c_1=0.05,\quad c_2=0.65,\quad\gamma=0.98,
\end{gathered}
\]
la política óptima usa producto 2 en los estados $s=5,\ldots,12$.
Las pruebas anteriores explican por qué tiene que existir una política
de intervalo; para encontrar sus extremos se evalúa la política y se
verifican las desigualdades de Bellman.

En el proyecto, el archivo
\begin{center}
\path{forgetting_random_replacement/code/verify_interval_theorem.py}
\end{center}
comprueba esta banda y las desigualdades de todas las acciones con
aritmética racional exacta. Por tanto, los signos del certificado no
dependen de una tolerancia numérica.

\subsection{Por qué el llenado inicial es diferente}

Si todavía solo se recuerdan $m<n$ observaciones, una compra lleva
a la capa de memoria $m+1$. El estado necesita ahora dos números:
$(m,s)$. La ventaja del producto 2 es
\[
G_{m,s}=-\Delta c+\gamma\Delta p
          \left(V_{m+1,s+1}-V_{m+1,s}\right).
\]
La ecuación ya no conecta solo vecinos dentro de una capa fija, como
en \eqref{eq:tridiagonal}. El resultado de intervalo puede fallar.

Por ejemplo, con
\[
\begin{gathered}
n=30,\quad p_0=0.45,\quad p_1=0.02,\quad p_2=0.15,\\
R=1,\quad c_1=0,\quad c_2=0.53,\quad\gamma=0.99985,
\end{gathered}
\]
en la capa $m=3$ las ventajas son aproximadamente
\[
G_{3,0}=0.03721,\quad G_{3,1}=-0.00979,\quad
G_{3,2}=-0.02710,\quad G_{3,3}=0.01320.
\]
Se elige producto 2 en $\{0,3\}$, que no es un intervalo. Los signos
de este ejemplo también están verificados con fracciones exactas en
el archivo indicado. Esta diferencia es la razón para formular ambas
pruebas específicamente para memoria llena.

\begin{thebibliography}{1}
\bibitem{ying-zhu2020}
Lexing Ying y Yuhua Zhu (2020).
\emph{A Note on Optimization Formulations of Markov Decision Processes}.
Sección 2: formulación primal y dual para problemas descontados.
\url{https://arxiv.org/abs/2012.09417}.
\end{thebibliography}

\end{document}
